\documentclass[11pt]{article}
\usepackage[a4paper,margin=1in]{geometry}
\usepackage{fontspec}
\usepackage[boldfont=false]{xeCJK}
\setCJKmainfont{FandolSong-Regular.otf}
\usepackage{amsmath}
\usepackage{amssymb}
\usepackage{array}
\usepackage{booktabs}
\usepackage{graphicx}
\usepackage{adjustbox}
\usepackage{enumitem}
\setlist[itemize]{topsep=2pt,partopsep=0pt,parsep=0pt,itemsep=2pt,leftmargin=1.4em}
\setlist[enumerate]{topsep=2pt,partopsep=0pt,parsep=0pt,itemsep=2pt,leftmargin=1.6em}
\usepackage{parskip}
\usepackage{fvextra}
\fvset{breaklines=true,breakanywhere=true,fontsize=\small}
\setlength{\parindent}{0pt}

\begin{document}

\begin{center}
  {\LARGE\bfseries Programming}\\[0.35em]
  {\large Igcse Computer Science}
\end{center}
\vspace{0.6em}

\section*{Variables and constants}

A \textbf{variable} 变量 is a named store that holds a value which \textbf{can change} while the program runs. A \textbf{constant} 常量 is a named store whose value is \textbf{fixed} and does not change.

You should \textbf{declare} 声明 them (say their name and type) before use:

\begin{Verbatim}
DECLARE score : INTEGER
DECLARE name : STRING
CONSTANT Pi = 3.142
\end{Verbatim}

Use a constant for a value that never changes (like Pi), so it is set in one place and is easy to read.

\section*{Data types}

A \textbf{data type} 数据类型 says what kind of value a variable holds. You must know five basic types.

\begin{center}
\adjustbox{max width=\linewidth}{%
\begin{tabular}{lll}
\toprule
\textbf{Type} & \textbf{Holds} & \textbf{Example} \\
\midrule
\textbf{integer} 整数 & a whole number & \texttt{42}, \texttt{-7} \\
\textbf{real} 实数 & a number with a decimal point & \texttt{3.14}, \texttt{-0.5} \\
\textbf{char} 字符 & a single character & \texttt{'A'}, \texttt{'?'} \\
\textbf{string} 字符串 & a sequence of characters & \texttt{"Hello"} \\
\textbf{Boolean} 布尔值 & one of two values & \texttt{TRUE} or \texttt{FALSE} \\
\bottomrule
\end{tabular}}
\end{center}

\section*{Input and output}

\textbf{Input} 输入 reads a value from the user. \textbf{Output} 输出 shows a value on the screen.

\begin{Verbatim}
OUTPUT "What is your name?"
INPUT name
OUTPUT "Hello ", name
\end{Verbatim}

\section*{The three basic structures}

Every program is built from three control structures.

\subsection*{Sequence}

\textbf{Sequence} 顺序 means the steps run one after another, in order, from top to bottom.

\begin{Verbatim}
INPUT length
INPUT width
area ← length * width
OUTPUT area
\end{Verbatim}

\subsection*{Selection}

\textbf{Selection} 选择 chooses which steps to run, based on a condition. Use an \texttt{IF} statement, or a \texttt{CASE} statement when there are many choices.

\begin{Verbatim}
IF score >= 50 THEN
    OUTPUT "Pass"
ELSE
    OUTPUT "Fail"
ENDIF
\end{Verbatim}

\begin{Verbatim}
CASE OF grade
    'A' : OUTPUT "Excellent"
    'B' : OUTPUT "Good"
    OTHERWISE : OUTPUT "Keep trying"
ENDCASE
\end{Verbatim}

\subsection*{Iteration}

\textbf{Iteration} 迭代 (a loop 循环) repeats steps. There are three kinds.

A \textbf{count-controlled loop} 计数循环 repeats a fixed number of times:

\begin{Verbatim}
FOR i ← 1 TO 5
    OUTPUT "Hello"
NEXT i
\end{Verbatim}

A \textbf{pre-condition loop} 前测循环 checks the condition \textbf{before} each repeat, so it may run zero times:

\begin{Verbatim}
WHILE answer <> "stop" DO
    INPUT answer
ENDWHILE
\end{Verbatim}

A \textbf{post-condition loop} 后测循环 checks the condition \textbf{after} each repeat, so it always runs \textbf{at least once}:

\begin{Verbatim}
REPEAT
    INPUT password
UNTIL password = "secret"
\end{Verbatim}

\section*{Totalling and counting}

\begin{itemize}
\item \textbf{totalling} 求和 — keep adding values to a running total: \texttt{total ← total + value}.

\item \textbf{counting} 计数 — add 1 to a counter each time: \texttt{count ← count + 1}.

\end{itemize}

\begin{Verbatim}
total ← 0
FOR i ← 1 TO 10
    INPUT mark
    total ← total + mark
NEXT i
OUTPUT total
\end{Verbatim}

\section*{Operators}

\subsection*{Arithmetic operators}

\begin{center}
\adjustbox{max width=\linewidth}{%
\begin{tabular}{ll}
\toprule
\textbf{Operator} & \textbf{Meaning} \\
\midrule
\texttt{+} \texttt{-} \texttt{*} \texttt{/} & add, subtract, multiply, divide \\
\texttt{\textasciicircum{}} & raised to the power of \\
\texttt{MOD} & the \textbf{remainder} 余数 after division \\
\texttt{DIV} & the whole-number part of a division \\
\bottomrule
\end{tabular}}
\end{center}

For example, \texttt{17 MOD 5} is \texttt{2}, and \texttt{17 DIV 5} is \texttt{3}.

\subsection*{Relational operators}

These compare two values and give a Boolean result: \texttt{=}, \texttt{\textless{}}, \texttt{\textless{}=}, \texttt{\textgreater{}}, \texttt{\textgreater{}=}, and \texttt{\textless{}\textgreater{}} (not equal to).

\subsection*{Logical operators}

These join conditions: \texttt{AND} (both must be true), \texttt{OR} (at least one must be true), \texttt{NOT} (reverses true/false).

\begin{Verbatim}
IF age >= 13 AND age <= 19 THEN
    OUTPUT "Teenager"
ENDIF
\end{Verbatim}

\section*{String handling}

A string is made of characters. Useful operations:

\begin{itemize}
\item \textbf{length} 长度 — the number of characters in the string;

\item \textbf{substring} 子串 — a smaller part taken from the string;

\item \textbf{upper case} 大写 — change letters to capitals;

\item \textbf{lower case} 小写 — change letters to small letters.

\end{itemize}

\begin{Verbatim}
name ← "Computer"
OUTPUT LENGTH(name)          // 8
OUTPUT SUBSTRING(name, 1, 4) // "Comp"
OUTPUT UCASE(name)           // "COMPUTER"
OUTPUT LCASE(name)           // "computer"
\end{Verbatim}

(The first character may be counted as position 0 or position 1, depending on the language.)

\section*{Nested statements}

A \textbf{nested statement} 嵌套语句 is one control structure placed \textbf{inside} another. You can nest selection and iteration.

\begin{Verbatim}
FOR i ← 1 TO 3
    IF i MOD 2 = 0 THEN
        OUTPUT i, " is even"
    ELSE
        OUTPUT i, " is odd"
    ENDIF
NEXT i
\end{Verbatim}

You will not have to write more than three levels of nesting.

\section*{Procedures and functions}

To avoid repeating code, you can break a program into named blocks.

\begin{itemize}
\item A \textbf{procedure} 过程 is a named block of code that does a task. You run it with \texttt{CALL}.

\item A \textbf{function} 函数 is like a procedure, but it \textbf{returns a value} back to where it was called.

\end{itemize}

A \textbf{parameter} 参数 is a value passed into a procedure or function (up to three parameters).

\begin{Verbatim}
PROCEDURE Greet(personName : STRING)
    OUTPUT "Hello ", personName
ENDPROCEDURE

CALL Greet("Sam")
\end{Verbatim}

\begin{Verbatim}
FUNCTION Square(n : INTEGER) RETURNS INTEGER
    RETURN n * n
ENDFUNCTION

answer ← Square(5)   // answer is 25
\end{Verbatim}

\subsection*{Local and global variables}

\begin{itemize}
\item A \textbf{local variable} 局部变量 is declared inside a procedure or function. It can only be used there.

\item A \textbf{global variable} 全局变量 is declared in the main program. It can be used anywhere.

\end{itemize}

Local variables are safer, because they cannot be changed by mistake from another part of the program.

\section*{Library routines}

A \textbf{library routine} 库程序 is a ready-made piece of code you can use. You must know these:

\begin{center}
\adjustbox{max width=\linewidth}{%
\begin{tabular}{ll}
\toprule
\textbf{Routine} & \textbf{What it does} \\
\midrule
\texttt{MOD} & gives the remainder of a division \\
\texttt{DIV} & gives the whole-number part of a division \\
\texttt{ROUND} & rounds a real number to a number of decimal places \\
\texttt{RANDOM} & gives a random number \\
\bottomrule
\end{tabular}}
\end{center}

\section*{Writing a maintainable program}

A \textbf{maintainable} 可维护的 program is easy for other people to read and change later. To make one:

\begin{itemize}
\item use \textbf{meaningful identifiers} 有意义的标识符 — clear names for variables, constants, arrays, procedures and functions (\texttt{totalScore}, not \texttt{x});

\item add \textbf{comments} 注释 to explain what parts of the code do;

\item use procedures and functions to split the work into small blocks.

\end{itemize}

\section*{Arrays}

An \textbf{array} 数组 is a single variable that holds \textbf{many values} of the same type, found by an \textbf{index} 索引 (a position number).

\subsection*{One-dimensional arrays}

A \textbf{one-dimensional (1D) array} 一维数组 is like a single list.

\begin{Verbatim}
DECLARE scores : ARRAY[1:5] OF INTEGER
scores[1] ← 90
scores[2] ← 75
OUTPUT scores[1]
\end{Verbatim}

You can use a loop to fill or read an array:

\begin{Verbatim}
FOR i ← 1 TO 5
    INPUT scores[i]
NEXT i
\end{Verbatim}

\subsection*{Two-dimensional arrays}

A \textbf{two-dimensional (2D) array} 二维数组 is like a table with rows and columns. It uses two indexes.

\begin{Verbatim}
DECLARE grid : ARRAY[1:3, 1:3] OF INTEGER
grid[1, 1] ← 5
grid[2, 3] ← 8
\end{Verbatim}

You read a 2D array with a loop inside a loop (nested iteration).

\section*{File handling}

A program can store data in a \textbf{file} 文件 so it is kept after the program stops. You must \textbf{open} the file, use it, then \textbf{close} it.

\begin{Verbatim}
OPENFILE "data.txt" FOR WRITE
WRITEFILE "data.txt", "Hello"
CLOSEFILE "data.txt"

OPENFILE "data.txt" FOR READ
READFILE "data.txt", line
CLOSEFILE "data.txt"
\end{Verbatim}

You can read and write single items of data or a whole line of text. Always close a file when you have finished with it.


\end{document}
